\documentclass[12pt]{article}
\usepackage[margin=1in]{geometry}
\usepackage{ctex} % 中文支持
\usepackage{tikz}
\usetikzlibrary{shapes.geometric, arrows.meta, positioning, fit, backgrounds}

\begin{document}

\begin{figure}[htbp]
\centering
% scale=0.85 是整体缩放比例，如果仍觉得大，可修改为 0.75 等更小数值
\begin{tikzpicture}[scale=0.85, transform shape, node distance=1cm and 1cm]

% 定义紧凑型方框样式 (使用 text width 强制文本换行，让方框变窄)
\tikzstyle{mainbox} =[rectangle, rounded corners, text width=3.2cm, minimum height=1.2cm, text centered, draw=black, thick, fill=blue!10, font=\small\bfseries]
\tikzstyle{subbox} =[rectangle, text width=2.8cm, minimum height=0.6cm, text centered, draw=gray, fill=white, font=\footnotesize]
\tikzstyle{arrow} =[thick, ->, >=stealth, draw=blue!70!black]
\tikzstyle{dashed_arrow} =[thick, dashed, ->, >=stealth, draw=red!70!black]

% 绘制主任务节点
\node (task1)[mainbox] {任务 1：数学模型\\ (理论框架)};
\node (task2)[mainbox, below left=1.8cm and -0.2cm of task1] {任务 2：敏感性分析\\ (实证验证)};
\node (task3)[mainbox, below right=1.8cm and -0.2cm of task1] {任务 3：消费者积分\\ (商业转化)};

% 绘制任务1的子节点
\node (t1_sub1)[subbox, above left=0.6cm and -1.2cm of task1] {避免与额外排放};
\node (t1_sub2)[subbox, above right=0.6cm and -1.2cm of task1] {因子与假设论证};
\draw [arrow] (t1_sub1) -- (task1);
\draw [arrow] (t1_sub2) -- (task1);

% 绘制任务2的子节点
\node (t2_sub1) [subbox, below=0.5cm of task2] {基准数据代入核算};
\node (t2_sub2)[subbox, below=0.2cm of t2_sub1] {控制变量扰动测试};
\draw [arrow] (task2) -- (t2_sub1);
\draw [arrow] (t2_sub1) -- (t2_sub2);

% 绘制任务3的子节点
\node (t3_sub1)[subbox, below=0.5cm of task3] {微观单件碳排映射};
\node (t3_sub2)[subbox, below=0.2cm of t3_sub1] {动态转化系数设计};
\draw [arrow] (task3) -- (t3_sub1);
\draw [arrow] (t3_sub1) -- (t3_sub2);

% 绘制任务之间的关联箭头与文字说明 (字号改为 scriptsize 避免拥挤)
\draw [arrow] (task1) -| node[anchor=east, align=center, font=\scriptsize, text=blue!80!black, yshift=0.3cm] {提供方程式\\与宏观参数} (task2);

\draw [arrow] (task1) -| node[anchor=west, align=center, font=\scriptsize, text=blue!80!black, yshift=0.3cm] {提供微观单品\\减排计算逻辑} (task3);

\draw [dashed_arrow] (task2.east) -- node[anchor=south, align=center, font=\scriptsize, text=red!80!black] {真实数据反馈影响权重} (task3.west);

\draw [dashed_arrow] (task2.north) -- node[anchor=east, align=center, font=\scriptsize, text=red!80!black] {验证鲁棒性} (task1.south -| task2.north);

% 绘制整体虚线背景框
\begin{scope}[on background layer]
    \node[rectangle, draw=gray, dashed, fill=gray!5, fit=(t1_sub1) (t1_sub2) (task1) (t2_sub2) (t3_sub2), inner sep=0.4cm] (bg) {};
    \node[anchor=north, font=\normalsize\bfseries, text=gray!80] at (bg.north) {临期折扣零售净碳排核算逻辑图};
\end{scope}

\end{tikzpicture}
\caption{任务一、二、三的逻辑传递与数据闭环关系}
\end{figure}

\end{document}